%% ----------------------------------------------------------------------------
%%  §6.2  I teoremi di Rolle e di Lagrange
%% ----------------------------------------------------------------------------
\section{I teoremi di Rolle e di Lagrange}
\label{sec:rolle-lagrange}

Il teorema di Fermat riguarda un singolo punto. I due teoremi di questo paragrafo
riguardano invece un intero intervallo, e sono la chiave di volta di tutto il capitolo:
mettono in relazione la variazione \emph{complessiva} $f(b)-f(a)$ di una funzione con
il valore della derivata in un punto opportuno. Il primo è un caso particolare del
secondo, ma è anche lo strumento con cui il secondo si dimostra.

\begin{teorema}[di Rolle]\label{teo:rolle}\index{teorema!di Rolle}
Sia $f$ una funzione continua nell'intervallo chiuso $[a,b]$ e derivabile
nell'intervallo aperto $(a,b)$. Se $f(a)=f(b)$, esiste un punto $x_{0}\in(a,b)$ tale
che
\begin{equation}\label{eq:rolle}
  f'(x_{0})=0 .
\end{equation}
\end{teorema}

\begin{schema}
\Ipotesi{\begin{itemize}[nosep, leftmargin=1.2em]
  \item $f$ continua in $[a,b]$;
  \item $f$ derivabile in $(a,b)$;
  \item $f(a)=f(b)$.
\end{itemize}}
\Tesi{esiste $x_{0}\in(a,b)$ con $f'(x_{0})=0$.}
\end{schema}

\begin{strategia}
Cerchiamo $x_{0}$ tra i punti di massimo e di minimo di $f$. La continuità
nell'intervallo chiuso e limitato ci permette di usare il teorema di Weierstrass: esistono
un punto di massimo e un punto di minimo assoluti. Se almeno uno dei due è
\emph{interno}, il teorema di Fermat (applicabile, perché $f$ è derivabile in $(a,b)$)
dà subito $f'(x_{0})=0$. Resta il caso in cui massimo e minimo cadono entrambi negli
estremi: qui interviene l'ipotesi $f(a)=f(b)$, che costringe massimo e minimo a
coincidere, cioè $f$ a essere costante; e una funzione costante ha derivata nulla
ovunque.
\end{strategia}

\begin{dimostrazione}
Poiché $f$ è continua in $[a,b]$, per il teorema di Weierstrass (\cref{teo:weierstrass})
esistono $x_{1},x_{2}\in[a,b]$ tali che
\begin{equation}\label{eq:rolle-weierstrass}
  m=f(x_{1})\le f(x)\le f(x_{2})=M \qquad\perogni x\in[a,b].
\end{equation}
Distinguiamo due casi.

\emph{Primo caso: almeno uno dei due punti $x_{1}$, $x_{2}$ è interno ad $[a,b]$.}
Supponiamo, per fissare le idee, che sia $x_{2}\in(a,b)$. Allora $x_{2}$ è un punto di
massimo (assoluto, e quindi relativo) interno, nel quale $f$ è derivabile; per il
teorema di Fermat (\cref{teo:fermat}) risulta $f'(x_{2})=0$, e la tesi vale con
$x_{0}=x_{2}$. Se invece è interno $x_{1}$, si ragiona allo stesso modo con il minimo.

\emph{Secondo caso: entrambi i punti $x_{1}$, $x_{2}$ sono estremi dell'intervallo.}
Allora $m$ e $M$ sono due dei valori $f(a)$, $f(b)$, che per ipotesi sono uguali;
quindi $m=M$, e dalla \eqref{eq:rolle-weierstrass} segue $f(x)=M$ per ogni $x\in[a,b]$.
La funzione è costante, e la sua derivata è nulla in ogni punto
\eqref{eq:derivata-costante}: la tesi vale con un qualunque $x_{0}\in(a,b)$.
\end{dimostrazione}

\begin{figure}[htb]
\centering
\begin{tikzpicture}
\begin{axis}[width=0.85\textwidth, height=5cm, xmin=0, xmax=6.4, ymin=0, ymax=3.3,
  axis lines=left, xtick={0.5,1.809,4.427,5.736},
  xticklabels={$a$,$x_{0}$,$x_{0}'$,$b$}, ytick=\empty, clip=false]
  \addplot[curva, domain=0.5:5.736, samples=150] {2+0.8*sin(deg(1.2*(x-0.5)))};
  \addplot[curva2, domain=1.0:2.6] {2.8};
  \addplot[curva2, domain=3.6:5.2] {1.2};
  \addplot[tratteggio] coordinates {(0.5,0) (0.5,2) (5.736,2) (5.736,0)};
  \addplot[tratteggio] coordinates {(1.809,0) (1.809,2.8)};
  \addplot[tratteggio] coordinates {(4.427,0) (4.427,1.2)};
  \addplot[only marks, mark=*, mark size=1.6pt, blu] coordinates
    {(0.5,2) (5.736,2) (1.809,2.8) (4.427,1.2)};
  \node[font=\scriptsize, grigio, anchor=west] at (axis cs:5.8,2.0) {$f(a)=f(b)$};
\end{axis}
\end{tikzpicture}
\caption{Il teorema di Rolle: se $f(a)=f(b)$, il grafico ha almeno un punto con
tangente orizzontale. Qui i punti sono due: il massimo $x_{0}$ e il minimo $x_{0}'$; il
teorema garantisce l'esistenza di almeno uno, non l'unicità.}
\label{fig:rolle}
\end{figure}

\begin{osservazione}[Le ipotesi del teorema di Rolle sono tutte necessarie]
\label{oss:rolle-ipotesi}
Se si toglie una qualunque delle tre ipotesi la tesi può essere falsa
(\cref{fig:rolle-ipotesi}).
\begin{enumerate}[label=\textup{(\alph*)}]
  \item \emph{Continuità in $[a,b]$.} La funzione $f(x)=x$ per $x\in[0,1)$, con
  $f(1)=0$, è derivabile in $(0,1)$, con derivata $1$, e verifica $f(0)=f(1)$; ma è
  discontinua in $x=1$, e la sua derivata non si annulla mai. La continuità serviva per
  il teorema di Weierstrass: infatti $f$ non ha massimo in $[0,1]$.
  \item \emph{Derivabilità in $(a,b)$.} La funzione $f(x)=1-\abs x$ in $[-1,1]$ è
  continua, verifica $f(-1)=f(1)=0$ e ha massimo nel punto interno $0$; ma lì non è
  derivabile, e altrove la derivata vale $\pm1$. Il teorema di Fermat non era
  applicabile.
  \item \emph{Uguaglianza $f(a)=f(b)$.} La funzione $f(x)=x$ in $[0,1]$ è continua e
  derivabile, ma $f'(x)=1$ per ogni $x$. Massimo e minimo cadono negli estremi senza che
  la funzione sia costante.
\end{enumerate}
Si noti invece che \emph{non} si richiede la derivabilità negli estremi: la funzione
$f(x)=\sqrt{1-x^{2}}$ in $[-1,1]$ (il cui grafico è una semicirconferenza) ha
derivata infinita in $\pm1$, ma verifica le ipotesi, e infatti $f'(0)=0$.
\end{osservazione}

\begin{figure}[htb]
\centering
\begin{tikzpicture}
\begin{axis}[name=r1, width=0.36\textwidth, height=0.32\textwidth, axis lines=middle,
  xmin=-0.3, xmax=1.35, ymin=-0.25, ymax=1.3, xtick={1}, ytick=\empty,
  title={\scriptsize (a) discontinua in $b$}]
  \addplot[curva, domain=0:1, samples=2] {x};
  \addplot[only marks, mark=o, mark size=1.6pt, blu, fill=white] coordinates {(1,1)};
  \addplot[only marks, mark=*, mark size=1.6pt, blu] coordinates {(1,0) (0,0)};
\end{axis}
\begin{axis}[at={(r1.east)}, anchor=west, xshift=0.6cm, name=r2, width=0.36\textwidth,
  height=0.32\textwidth, axis lines=middle, xmin=-1.35, xmax=1.35, ymin=-0.25,
  ymax=1.3, xtick={-1,1}, ytick=\empty, title={\scriptsize (b) non derivabile in $0$}]
  \addplot[curva, domain=-1:1, samples=3] {1-abs(x)};
\end{axis}
\begin{axis}[at={(r2.east)}, anchor=west, xshift=0.6cm, width=0.36\textwidth,
  height=0.32\textwidth, axis lines=middle, xmin=-0.3, xmax=1.35, ymin=-0.25,
  ymax=1.3, xtick={1}, ytick=\empty, title={\scriptsize (c) $f(a)\neq f(b)$}]
  \addplot[curva, domain=0:1, samples=2] {x};
  \addplot[only marks, mark=*, mark size=1.6pt, blu] coordinates {(0,0) (1,1)};
\end{axis}
\end{tikzpicture}
\caption{Tre funzioni che violano una sola delle ipotesi del teorema di Rolle e per le
quali la tesi è falsa: nessuna ha un punto a tangente orizzontale.}
\label{fig:rolle-ipotesi}
\end{figure}

Il teorema di Rolle ha un'interpretazione geometrica evidente: se la corda che
congiunge gli estremi del grafico è orizzontale, in qualche punto anche la tangente lo
è. Se ora «inclinassimo» la figura, ci aspetteremmo di trovare un punto in cui la
tangente è \emph{parallela alla corda}. È il contenuto del teorema seguente.

\begin{teorema}[di Lagrange, o del valor medio]\label{teo:lagrange}
\index{teorema!di Lagrange}\index{teorema!del valor medio}
Sia $f$ una funzione continua nell'intervallo chiuso $[a,b]$ e derivabile
nell'intervallo aperto $(a,b)$. Esiste un punto $x_{0}\in(a,b)$ tale che
\begin{equation}\label{eq:lagrange}
  \frac{f(b)-f(a)}{b-a}=f'(x_{0}) .
\end{equation}
\end{teorema}

\begin{schema}
\Ipotesi{$f$ continua in $[a,b]$ e derivabile in $(a,b)$.}
\Tesi{esiste $x_{0}\in(a,b)$ con $f(b)-f(a)=f'(x_{0})\,(b-a)$.}
\end{schema}

\begin{strategia}
Il primo membro della \eqref{eq:lagrange} è la pendenza della corda che congiunge
$(a,f(a))$ e $(b,f(b))$; il secondo è la pendenza della tangente in $x_{0}$. Ci
riconduciamo al teorema di Rolle \emph{sottraendo la corda}: la funzione ausiliaria
$g(x)=f(x)-r(x)$, dove $y=r(x)$ è l'equazione della retta secante, misura la distanza
verticale tra il grafico e la corda. Questa distanza è nulla negli estremi, quindi $g$
verifica le ipotesi del teorema di Rolle; nel punto in cui $g'$ si annulla, la
pendenza di $f$ eguaglia quella della retta. È un procedimento tipico: per dimostrare
un risultato più generale, si costruisce una funzione alla quale si può applicare il
caso particolare già noto.
\end{strategia}

\begin{dimostrazione}
La retta passante per i punti $(a,f(a))$ e $(b,f(b))$ ha equazione
$y=r(x)=f(a)+\frac{f(b)-f(a)}{b-a}(x-a)$. Consideriamo la funzione
\begin{equation}\label{eq:lagrange-g}
  g(x)=f(x)-\Bigl[f(a)+\frac{f(b)-f(a)}{b-a}(x-a)\Bigr], \qquad x\in[a,b].
\end{equation}
La funzione $g$ è continua in $[a,b]$ e derivabile in $(a,b)$, perché differenza di $f$
e di un polinomio di primo grado (\cref{teo:operazioni-continue,teo:operazioni-derivate});
inoltre
\[
  g(a)=f(a)-f(a)=0, \qquad g(b)=f(b)-\bigl[f(a)+f(b)-f(a)\bigr]=0 .
\]
Per il teorema di Rolle (\cref{teo:rolle}) esiste $x_{0}\in(a,b)$ tale che
$g'(x_{0})=0$. Poiché
\[
  g'(x)=f'(x)-\frac{f(b)-f(a)}{b-a} \qquad\perogni x\in(a,b),
\]
la condizione $g'(x_{0})=0$ è esattamente la \eqref{eq:lagrange}.
\end{dimostrazione}

\begin{figure}[htb]
\centering
\begin{tikzpicture}
\begin{axis}[width=0.85\textwidth, height=5.4cm, xmin=0, xmax=4.8, ymin=0, ymax=4.6,
  axis lines=left, xtick={0.25,1.5625,4}, xticklabels={$a$,$x_{0}$,$b$},
  ytick=\empty, clip=false]
  \addplot[curva, domain=0.1:4.3, samples=120] {2*sqrt(x)};
  \addplot[curva2, domain=0.25:4] {1+0.8*(x-0.25)};
  \addplot[curva3, domain=0.2:3.6] {2.5+0.8*(x-1.5625)};
  \addplot[tratteggio] coordinates {(0.25,0) (0.25,1)};
  \addplot[tratteggio] coordinates {(4,0) (4,4)};
  \addplot[tratteggio] coordinates {(1.5625,0) (1.5625,2.5)};
  \addplot[only marks, mark=*, mark size=1.6pt, blu] coordinates
    {(0.25,1) (4,4) (1.5625,2.5)};
  \node[font=\scriptsize, rame, anchor=north west] at (axis cs:2.5,2.75) {corda};
  \node[font=\scriptsize, verde!50!black, anchor=south east] at (axis cs:3.2,3.95)
    {tangente parallela};
\end{axis}
\end{tikzpicture}
\caption{Il teorema di Lagrange per $f(x)=2\sqrt x$ in $[a,b]=[\frac14,4]$: la corda ha
pendenza $\frac{4-1}{4-1/4}=\frac45$, e la tangente nel punto $x_{0}=\frac{25}{16}$, dove
$f'(x_{0})=\frac1{\sqrt{x_{0}}}=\frac45$, le è parallela.}
\label{fig:lagrange}
\end{figure}

\begin{idea}[Il «Tutor» e il velivolo: velocità media e velocità istantanea]
Se $s(t)$ è la legge oraria di un moto rettilineo, il primo membro della
\eqref{eq:lagrange} è la velocità media nell'intervallo di tempo $[a,b]$, il secondo la
velocità istantanea all'istante $x_{0}$. Il teorema di Lagrange afferma dunque che
\emph{esiste un istante in cui la velocità istantanea è uguale alla velocità media}.
È il principio su cui si fonda il sistema «Tutor» che controlla la velocità sulle
autostrade italiane: il sistema misura soltanto i tempi di passaggio sotto due portali
distanti, ad esempio, $20$~km; se l'automobile li ha percorsi in $8$~min, la
velocità media è di $150$~km/h e, per il teorema di Lagrange, il tachimetro ha
segnato $150$~km/h in almeno un istante, anche se nessuno l'ha visto. Allo stesso
modo, un velivolo che percorre i $300$~km tra due punti di riporto in un'ora ha
avuto, almeno in un istante, una velocità al suolo (lungo la rotta) di esattamente
$300$~km/h.
\end{idea}

Il punto $x_{0}$ di cui il teorema garantisce l'esistenza in generale non si può
calcolare esplicitamente, ma spesso si può trovare numericamente risolvendo
l'equazione $f'(x)=\frac{f(b)-f(a)}{b-a}$, ad esempio con il metodo di bisezione del
\cref{sec:bisezione} o con la funzione \texttt{brentq} di \texttt{scipy}, purché $f'$
sia continua e la differenza $f'(x)-\frac{f(b)-f(a)}{b-a}$ abbia segni opposti negli
estremi (per $x^{3}$ in $[-1,1]$, ad esempio, ciò non accade, benché i punti
$\pm\frac1{\sqrt3}$ esistano). Il programma
seguente lo fa per $f(x)=x^{3}$ in $[0,2]$, dove l'equazione $3x^{2}=4$ dà
$x_{0}=\frac2{\sqrt3}$, e per $f(x)=\sen x$ in $[0,\frac\pi2]$, dove $\cos x_{0}=\frac2\pi$.

\begin{python}[Il punto di Lagrange calcolato numericamente]
import math
from scipy.optimize import brentq

def xi_lagrange(f, df, a, b):
    """Un punto xi in (a, b) con f'(xi) = [f(b) - f(a)] / (b - a)."""
    m = (f(b) - f(a)) / (b - a)          # pendenza della corda
    return brentq(lambda t: df(t) - m, a + 1e-12, b - 1e-12), m

xi, m = xi_lagrange(lambda t: t**3, lambda t: 3 * t**2, 0, 2)
print(f"x^3:   corda = {m:.4f}   xi = {xi:.6f}"
      f"   2/sqrt(3)    = {2 / math.sqrt(3):.6f}")
xi, m = xi_lagrange(math.sin, math.cos, 0, math.pi / 2)
print(f"sen x: corda = {m:.4f}   xi = {xi:.6f}"
      f"   arccos(2/pi) = {math.acos(2 / math.pi):.6f}")

# stima di sqrt(101) con il teorema di Lagrange (vedi il testo)
sup = 10 + 1 / 20
inf = 10 + 1 / (2 * sup)
print(f"{inf:.6f} < sqrt(101) = {math.sqrt(101):.6f} < {sup:.6f}")
\end{python}
\begin{risultato}
x^3:   corda = 4.0000   xi = 1.154701   2/sqrt(3)    = 1.154701
sen x: corda = 0.6366   xi = 0.880689   arccos(2/pi) = 0.880689
10.049751 < sqrt(101) = 10.049876 < 10.050000
\end{risultato}

Posto $a=x_{0}$, $b=x_{0}+h$ (con $h>0$) e scritto il punto intermedio nella forma
$x_{0}+\vartheta h$, con $0<\vartheta<1$, la tesi del teorema di Lagrange diventa la
\emph{formula degli accrescimenti finiti}: esiste $\vartheta\in(0,1)$ tale che
\begin{equation}\label{eq:accrescimenti-finiti}
  f(x_{0}+h)=f(x_{0})+f'(x_{0}+\vartheta h)\,h ,
\end{equation}
valida, con lo stesso ragionamento applicato all'intervallo $[x_{0}+h,x_{0}]$, anche
per $h<0$: in entrambi i casi basta che $f$ sia continua nell'intervallo chiuso di
estremi $x_{0}$ e $x_{0}+h$ e derivabile nei suoi punti interni. Il confronto con
l'approssimazione lineare del \cref{sec:retta-tangente},
$f(x_{0}+h)=f(x_{0})+f'(x_{0})h+o(h)$, è istruttivo: lì la derivata è calcolata nel
punto noto $x_{0}$ e si commette un errore $o(h)$; qui l'uguaglianza è esatta, ma la
derivata è calcolata in un punto sconosciuto. Proprio questo la rende adatta a
ottenere \emph{stime}, come nel corollario seguente.

\needspace{9\baselineskip}
\begin{corollario}[Stima degli incrementi]\label{cor:lipschitz}
Sia $f$ derivabile in un intervallo $I$ e supponiamo che esista $M\ge0$ tale che
$\abs{f'(x)}\le M$ per ogni $x\in I$. Allora
\begin{equation}\label{eq:lipschitz}
  \abs{f(x)-f(y)}\le M\abs{x-y} \qquad\perogni x,y\in I .
\end{equation}
\end{corollario}

\begin{schema}
\Ipotesi{$f$ derivabile nell'intervallo $I$; \ $\abs{f'}\le M$ in $I$.}
\Tesi{$\abs{f(x)-f(y)}\le M\abs{x-y}$ per ogni $x,y\in I$.}
\end{schema}

\begin{strategia}
Basta applicare il teorema di Lagrange all'intervallo di estremi $x$ e $y$ (contenuto in
$I$, perché $I$ è un intervallo) e maggiorare la derivata nel punto intermedio.
\end{strategia}

\begin{dimostrazione}
Se $x=y$ non c'è nulla da dimostrare. Se, ad esempio, $x<y$, l'intervallo $[x,y]$ è
contenuto in $I$, e in esso $f$ è derivabile, quindi anche continua
(\cref{teo:derivabile-continua}). Per il teorema di Lagrange esiste $x_{0}\in(x,y)$ con
$f(y)-f(x)=f'(x_{0})(y-x)$, e quindi
$\abs{f(y)-f(x)}=\abs{f'(x_{0})}\,\abs{y-x}\le M\abs{y-x}$.
\end{dimostrazione}

Una funzione che verifica la \eqref{eq:lipschitz} per qualche costante $M$ si dice
\emph{lipschitziana}\index{funzione!lipschitziana} in $I$, dal nome del matematico tedesco
Rudolf Lipschitz. Poiché $\abs{D\sen x}=\abs{\cos x}\le1$ e
$\abs{D\operatorname{arctg}x}=\frac1{1+x^{2}}\le1$, si ottengono le disuguaglianze, valide
per ogni $x,y\in\R$,
\begin{equation}\label{eq:sen-arctg-lipschitz}
  \abs{\sen x-\sen y}\le\abs{x-y}, \qquad
  \abs{\operatorname{arctg}x-\operatorname{arctg}y}\le\abs{x-y} ,
\end{equation}
la prima delle quali era stata ricavata per via trigonometrica nell'\cref{es:seno-lipschitz}.
In termini applicativi il corollario è una regola di \emph{propagazione degli errori}:
se una grandezza $x$ è misurata con un errore non superiore a $\varepsilon$, il valore
calcolato $f(x)$ è affetto da un errore non superiore a $M\varepsilon$, dove $M$ è il
massimo di $\abs{f'}$ nell'intervallo di incertezza.

\begin{esempio}[Una stima di $\sqrt{101}$]\label{es:stima-radice}
Applichiamo il teorema di Lagrange a $f(x)=\sqrt x$ nell'intervallo $[100,101]$:
esiste $x_{0}\in(100,101)$ tale che
\[
  \sqrt{101}-10=\frac{1}{2\sqrt{x_{0}}}\,(101-100)=\frac{1}{2\sqrt{x_{0}}} .
\]
Poiché $\sqrt{x_{0}}>10$, risulta $\sqrt{101}-10<\frac1{20}$, cioè $\sqrt{101}<10{,}05$.
Ma allora anche $\sqrt{x_{0}}<\sqrt{101}<10{,}05$, e quindi
$\sqrt{101}-10>\frac{1}{20{,}1}$. In conclusione
$10{,}04975<\sqrt{101}<10{,}05$, come conferma l'ultima riga dell'output precedente.
Si noti come la stima superiore sia stata usata per migliorare quella inferiore.
\end{esempio}

\needspace{8\baselineskip}
\begin{notastorica}[Rolle, Lagrange, Bonnet]
Michel Rolle (1652--1719), membro dell'Académie des Sciences, enunciò nel 1691 il
teorema che porta il suo nome soltanto per i \emph{polinomi}, nel contesto della
ricerca delle radici delle equazioni algebriche: tra due radici consecutive di un
polinomio cade una radice di un polinomio da esso derivato secondo regole puramente
algebriche (la «cascata» di Rolle). Paradossalmente, Rolle fu per anni uno dei più
accesi critici del nuovo calcolo infinitesimale, che giudicava privo di fondamenti
solidi. Il nome «teorema di Rolle» fu attribuito soltanto nell'Ottocento. Il teorema
del valor medio compare nella \emph{Théorie des fonctions analytiques} (1797) di
Joseph-Louis Lagrange, in una forma legata allo sviluppo in serie delle funzioni;
Cauchy, nel 1823, ne diede una dimostrazione che presupponeva la continuità della
derivata. La dimostrazione che abbiamo presentato, che deduce il teorema da quello di
Rolle mediante la funzione ausiliaria \eqref{eq:lagrange-g}, è dovuta a Ossian Bonnet e
fu pubblicata nel \emph{Cours de calcul différentiel et intégral} di Joseph Serret
(1868); da allora è quella di tutti i manuali.
\end{notastorica}
